% Source: physical PDF pages 59--64; printed pages 36--41.
% Section 15.8 begins on physical p. 59 and continues through the end of physical p. 64.
% Figures: none. Source uncertainties: none.
\subsection[Generalizations of BRST Symmetry]{Generalizations of BRST Symmetry\footnotemark}
\footnotetext{This section lies somewhat out of the book's main line of development, and may be omitted in a first reading.}

The BRST symmetry described in the previous section has a useful generalization to the quantization of a wide class of theories, including general relativity and string theories. In all these cases, we deal with an action $I[\phi]$ and measure $[d\phi]\equiv\prod_r \dd\phi^r$ that are invariant under a set of infinitesimal transformations
\begin{equation}
 \phi^r\longrightarrow\phi^r+\epsilon^A\delta_A\phi^r.
 \tag{15.8.1}\label{eq:15.8.1}
\end{equation}

This is an abbreviated `De Witt' notation, with $r$ and $A$ including spacetime coordinates as well as discrete labels, and sums including integrals over these coordinates. For instance, for the gauge transformation \eqref{eq:15.1.9}, the index $A$ consists of a group index $a$ and a spacetime coordinate $x$, with $\epsilon^{a x}\equiv\epsilon^a(x)$, while the index $r$ consists of a vector index $\mu$ as well as a group index $a$ and a spacetime coordinate $x$, with $\phi^{\mu a x}\equiv A_\mu^a(x)$; in the notation of Eq.~\eqref{eq:15.8.1}, the variation $\delta_A\phi^r$ in the transformation \eqref{eq:15.1.9} reads
\[
 \delta_{b y}\phi^{\mu a x}
 =\delta_a^b\frac{\partial}{\partial x^\mu}\delta^4(x-y)
 +C^b{}_{ca}\phi^{\mu c x}\delta^4(x-y).
\]
As in the special case of Yang--Mills theories discussed in the previous section, BRST invariance can be used as a substitute for the Faddeev--Popov--De Witt formulation of these theories, one that is applicable even where the Faddeev--Popov--De Witt approach fails. Nevertheless, in order to motivate the introduction of BRST invariance, we shall begin here with the Faddeev--Popov--De Witt formulation of theories with general local symmetries, and then go on to consider further generalizations.

By following the same arguments used to derive Eq.~\eqref{eq:15.5.21}, we obtain the general Faddeev--Popov--De Witt theorem:
\begin{equation}
 \frac{C}{\Omega}\int[d\phi]e^{\ii I[\phi]}V[\phi]
 =\int[d\phi]e^{\ii I[\phi]}B\{f[\phi]\}\operatorname{Det}(\delta_Af_B[\phi])V[\phi],
 \tag{15.8.2}\label{eq:15.8.2}
\end{equation}
where $V[\phi]$ is an arbitrary functional of $\phi^r$ that is invariant under the gauge transformations \eqref{eq:15.8.1}; $f_A[\phi]$ are a set of gauge-fixing functionals\footnote{We use the same letters $A$, $B$, etc. to label the $f_A$ as the gauge variations $\delta_A$, in order to emphasize that there must be as many gauge-fixing functionals as there are independent gauge transformations. However, in some cases like string theory it is natural to use gauge-fixing functionals $f^a$ for which, although the index $a$ runs over `as many' values as the index $A$ on the gauge variations $\delta_A$, the values taken by these indices are quite different. No change is needed in the present formalism as long as we can define $f_A=c_{Aa}f^a$, with $c_{Aa}$ field-independent and non-singular.} of the $\phi^r$, chosen so that the `matrix' $\delta_Af_B[\phi]$ has a non-vanishing determinant, and $B[f]$ is a more-or-less arbitrary functional of the $f_A$ (as for instance $\prod_A\delta(f_A)$). The constant $\Omega$ is the volume of the gauge group, and the constant $C$ is defined (as in Eq.~\eqref{eq:15.5.19}) by
\begin{equation}
 C\equiv\int[df]B[f].
 \tag{15.8.3}\label{eq:15.8.3}
\end{equation}
As we have seen, in gauge theories the importance of Eq.~\eqref{eq:15.8.2} is that it tells us that the integral on the right-hand side is independent of the choice of the gauge-fixing functionals $f_A$, and depends on $B[f]$ only through the constant $C$. Where some meaning can be given to the usually infinite group volume $\Omega$, as in gauge theories on a finite spacetime lattice, Eq.~\eqref{eq:15.8.2} can also have value as a formula for the integral on the left-hand side.

To define a nilpotent BRST transformation, we must first express the functional $B[f]$ as a Fourier transform
\begin{equation}
 B[f]=\int[dh]\exp(ih^Af_A)\mathcal B[h],
 \tag{15.8.4}\label{eq:15.8.4}
\end{equation}
where $[dh]\equiv\prod_A \dd h^A$. Also, the determinant can be expressed as an integral over fermionic c-number fields\footnote{It is common in string theories and elsewhere to find the ghost fields $\omega^{*A}$ and $\omega^A$ written $b^A$ (or $b_A$) and $c^A$, respectively.} $\omega^{*A}$ and $\omega^A$:
\begin{equation}
 \operatorname{Det}(\delta_Af_B[\phi])\propto
 \int[d\omega^*][d\omega]\exp\left(\ii\omega^{*B}\omega^A\delta_Af_B\right),
 \tag{15.8.5}\label{eq:15.8.5}
\end{equation}
where $[d\omega^*]\equiv\prod_A \dd\omega^{*A}$ and $[d\omega]\equiv\prod_A \dd\omega^A$, and as usual `$\propto$' means proportional with field-independent factors. Inserting these in Eq.~\eqref{eq:15.8.2} gives a general formula for the gauge-fixed path integral
\begin{align}
 &\int[d\phi]\exp\bigl(\ii I[\phi]\bigr)B\{f[\phi]\}\operatorname{Det}(\delta_Af_B[\phi])V[\phi]\notag\\
 &\qquad\propto\int[d\phi][dh][d\omega^*][d\omega]
 \exp\bigl(\ii I_{\mathrm{NEW}}[\phi,h,\omega,\omega^*]\bigr)\mathcal B[h]V[\phi],
 \tag{15.8.6}\label{eq:15.8.6}
\end{align}
where $I_{\mathrm{NEW}}$ is the new total action:
\begin{equation}
 I_{\mathrm{NEW}}[\phi,h,\omega,\omega^*]
 =I[\phi]+h^Af_A[\phi]+\omega^{*B}\omega^A\delta_Af_B[\phi].
 \tag{15.8.7}\label{eq:15.8.7}
\end{equation}

As mentioned in Section 15.6, we can think of the ghost fields as compensation for the fact that we are integrating over all $\phi^r$, including those $\phi^r$ that differ only by gauge transformations \eqref{eq:15.8.1}. Because ghosts are fermions, loops of ghost lines carry extra minus signs that allow these loops to compensate for the integration over gauge-equivalent $\phi$s. But for this to work, there must be just as many ghost fields $\omega^A$ as there are independent gauge transformations. That is, since the $\omega^A$ are independent, the gauge transformations \eqref{eq:15.8.1} must all be independent. This is the case for gauge transformations in Yang--Mills theory and coordinate transformations in general relativity, but not always. The classic example of a theory with non-independent gauge transformations is the theory of $p$-form gauge fields, described in Section 8.8. A $p$-form $A$ (an antisymmetric tensor of rank $p$) undergoes a gauge transformation $A\longrightarrow A+\dd\phi$, where $\phi$ is a $(p-1)$-form, and $\dd\phi$ is its exterior derivative (the antisymmetrized derivative). Because $d$ is nilpotent, for $p\geq2$ we can shift $\phi$ by an amount $\dd\psi$ without changing the gauge transformation, so there is a sort of invariance under gauge transformations of gauge transformations, in which the transformation parameters are the $(p-2)$-forms $\psi$. In such cases we must compensate for introducing too many ghosts by also introducing `ghosts of ghosts'~\hyperref[ch15-ref-15]{[15]}. For $p\geq3$ we need to compensate further by introducing `ghosts of ghosts of ghosts,' and so on. In what follows we shall assume that the gauge transformations \eqref{eq:15.8.1} are all independent, so that the ghost fields $\omega^A$ (and the antighosts $\omega^{*A}$) are all we need.

Although the original symmetry \eqref{eq:15.8.1} has been eliminated by the insertion of the non-gauge-invariant functional $B[f]$, the new total action has an exact symmetry under the infinitesimal BRST transformations
\begin{equation}
 \chi\longrightarrow\chi+\theta s\chi,
 \tag{15.8.8}\label{eq:15.8.8}
\end{equation}
where $\chi$ is any of the $\phi^r$, $\omega^A$, $\omega^{*A}$, or $h^A$; $\theta$ is an infinitesimal anticommuting c-number; and $s$ is the \emph{Slavnov operator}
\begin{equation}
 s=\omega^A\delta_A\phi^r\frac{\delta_L}{\delta\phi^r}
 -\frac{1}{2}\omega^B\omega^Cf^A{}_{BC}\frac{\delta_L}{\delta\omega^A}
 -h^A\frac{\delta_L}{\delta\omega^{*A}}.
 \tag{15.8.9}\label{eq:15.8.9}
\end{equation}
In Eq.~\eqref{eq:15.8.9} the subscript $L$ denotes left differentiation, defined so that if $\delta F=\delta\chi G$, then $\delta_LF/\delta\chi=G$, and $f^A{}_{BC}$ is the structure constant\footnote{For instance, for a gauge transformation acting on a matter field $\psi(x)$ we have $\delta_{b y}\psi(x)=\ii t_b\psi(x)\delta^4(x-y)$, and so $\delta_{a z}\delta_{b y}\psi(x)=-t_a t_b\psi(x)\delta^4(x-y)\delta^4(x-z)$. Hence in this case we have $f^{c x}{}_{b y,a z}=C^c{}_{ba}\delta^4(x-y)\delta^4(x-z)$.} appearing in the commutation relation
\begin{equation}
 [\delta_B,\delta_C]=f^A{}_{BC}\delta_A.
 \tag{15.8.10}\label{eq:15.8.10}
\end{equation}
The $f^A{}_{BC}$ are field-independent in non-Abelian gauge theories and in string theories, though not always, but the BRST formalism is not limited to this case. A straightforward calculation gives
\begin{align}
 s^2={}&\frac{1}{2}\omega^A\omega^B
 \left[\delta_A\phi^s\frac{\delta_L(\delta_B\phi^r)}{\delta\phi^s}
 -\delta_B\phi^s\frac{\delta_L(\delta_A\phi^r)}{\delta\phi^s}
 -f^C{}_{AB}\delta_C\phi^r\right]\frac{\delta_L}{\delta\phi^r}\notag\\
 &-\frac{1}{2}\omega^B\omega^C\omega^D
 \left[f^E{}_{BC}f^A{}_{DE}
 +\delta_D\phi^r\frac{\delta_Lf^A{}_{BC}}{\delta\phi^r}\right]
 \frac{\delta_L}{\delta\omega^A}.
 \tag{15.8.11}\label{eq:15.8.11}
\end{align}
Hence the condition that the BRST transformation be nilpotent is equivalent to the commutation relation \eqref{eq:15.8.10}, together with a consistency condition
\begin{equation}
 f^E{}_{[BC}f^A{}_{D]E}
 +\delta_{[D}\phi^r\left(\delta_Lf^A{}_{BC]}/\delta\phi^r\right)=0,
 \tag{15.8.12}\label{eq:15.8.12}
\end{equation}
where the brackets in subscripts indicate antisymmetrization with respect to the enclosed indices $B$, $C$, and $D$. Eq.~\eqref{eq:15.8.12} may be derived from the commutation relation \eqref{eq:15.8.10} in the same way as the usual Jacobi identity, and takes the place of the Jacobi identity for symmetries with field-dependent structure constants.

To show that the transformation \eqref{eq:15.8.8} is a symmetry of $I_{\mathrm{NEW}}$, we note (recalling that $\theta$ anticommutes with $\omega^{*A}$) that Eq.~\eqref{eq:15.8.7} may be rewritten
\begin{equation}
 I_{\mathrm{NEW}}[\phi,h,\omega,\omega^*]=I[\phi]-s(\omega^{*A}f_A).
 \tag{15.8.13}\label{eq:15.8.13}
\end{equation}
The term $I[\phi]$ is BRST-invariant, because on the fields $\phi^r$ a BRST transformation is just a gauge transformation \eqref{eq:15.8.1} with $\epsilon^A$ replaced with $\theta\omega^A$, which commutes with all $\phi^r$. The term $s(\omega^{*A}f_A)$ is BRST-invariant because BRST transformations are nilpotent.

For several reasons we may need to consider a wider class of actions than those that can be constructed by the Faddeev--Popov--De Witt approach, by simply requiring that the action is invariant under the BRST transformation \eqref{eq:15.8.8}. As a step toward showing that such an action yields physically sensible results, we shall now prove the general result (already used in the previous section) that the most general BRST-invariant functional of ghost number zero is the sum of a functional of the $\phi$ alone, plus a term given by the action of the BRST operator $s$ on an arbitrary functional $\Psi$ of ghost number $-1$:
\begin{equation}
 I_{\mathrm{NEW}}[\phi,\omega,\omega^*,h]
 =I_0[\phi]+s\Psi[\phi,\omega,\omega^*,h]
 \tag{15.8.14}\label{eq:15.8.14}
\end{equation}
as for instance in the Faddeev--Popov--De Witt action \eqref{eq:15.8.13}. In brief, the BRST cohomology consists of gauge-invariant functionals $I[\phi]$ of the fields $\phi^r$ alone.

To prove Eq.~\eqref{eq:15.8.14}, we note that the BRST transformation \eqref{eq:15.8.8}--\eqref{eq:15.8.9} does not change the total number of $h^A$ and $\omega^{*A}$ fields, so if we expand $I$ in a series of terms $I_N$ that contain definite total numbers $N$ of $h^A$ and $\omega^{*A}$ fields, then there can be no cancellations in $sI$ between terms with different $N$, so each term must be separately BRST-invariant:
\begin{equation}
 sI_N=0.
 \tag{15.8.15}\label{eq:15.8.15}
\end{equation}
We next introduce what is called a Hodge operator:
\begin{equation}
 t\equiv\omega^{*A}\frac{\delta}{\delta h^A}.
 \tag{15.8.16}\label{eq:15.8.16}
\end{equation}
It is straightforward to check the anticommutation relation
\begin{equation}
 \{s,t\}=-\omega^{*A}\frac{\delta_L}{\delta\omega^{*A}}-h^A\frac{\delta}{\delta h^A}.
 \tag{15.8.17}\label{eq:15.8.17}
\end{equation}
Applying the operator $\{s,t\}$ to $I_N$ and using Eq.~\eqref{eq:15.8.15} then gives
\begin{equation}
 stI_N=-NI_N,
 \tag{15.8.18}\label{eq:15.8.18}
\end{equation}
so each $I_N$ except for $I_0$ is BRST-exact, in the sense that it may be written as the operator $s$ acting on some other functional. The complete functional $I$ may therefore be written in the form $I_0+s\Psi$, with
\begin{equation}
 \Psi=-\sum_{N=1}^{\infty}\frac{tI_N}{N}.
 \tag{15.8.19}\label{eq:15.8.19}
\end{equation}
The term $I_0$ is by definition independent of $\omega^{*A}$ and $h^A$, and since we assume it has zero ghost number it must also be independent of $\omega^A$, as was to be proved.

To show the invariance of physical matrix elements under changes in the definition of the gauge-fixing functional $\Psi$, we define a fermionic `charge' $Q$, such that the change under a BRST transformation of any operator $\Phi$ is
\begin{equation}
 \delta_\theta\Phi=\ii[\theta Q,\Phi]=\ii\theta[Q,\Phi]_{\mp}
 \tag{15.8.20}\label{eq:15.8.20}
\end{equation}
with the top or bottom sign in $[x,y]_{\mp}\equiv xy\mp yx$ according as $\Phi$ is bosonic or fermionic. Just as in the previous section, the nilpotence of the BRST transformation then tells us that $Q^2=0$. Matrix elements of gauge-invariant operators between physical states will be independent of the choice of $\Psi$ if and only if the physical states $\ket{\alpha}$ and $\bra{\beta}$ satisfy
\begin{equation}
 Q\ket{\alpha}=\bra{\beta}Q=0,
 \tag{15.8.21}\label{eq:15.8.21}
\end{equation}
so that physically distinguishable physical states are again in one-to-one correspondence with elements of the cohomology of $Q$. The general BRST-invariant action \eqref{eq:15.8.3} will therefore yield physically sensible results for any gauge-fixing functional $\Psi$ if we can find \emph{some} $\Psi$, like axial gauge in Yang--Mills theories, in which the ghosts do not interact with other fields. If this ghost-free choice of $\Psi$ is inconvenient for actual calculation, as for instance axial gauge is inconvenient because it violates Lorentz invariance, we can adopt any gauge-fixing $\Psi$ we like, and still be confident that there is a unitary S-matrix with no ghosts in initial or final states.

This approach works well in string theories, where so-called light-cone quantization takes the place of axial gauge. But in other theories like general relativity there is no way of choosing a coordinate system in which the ghosts decouple. Such theories may be dealt with by the BRST-quantization method described at the end of the previous section, using BRST invariance to prove that the S-matrix in a physical ghost-free Hilbert space is unitary.

The discovery~\hyperref[ch15-ref-17]{[17]} of invariance under an `anti-BRST' symmetry~\hyperref[ch15-ref-18]{[18]} showed that, despite appearances, there is a similarity between the roles of $\omega^A$ and $\omega^{*A}$, which remains somewhat mysterious.
